\documentclass[12pt]{article}
\usepackage{amsmath,amssymb,amsthm,hyperref,booktabs,xcolor,enumitem}
\usepackage[margin=1in]{geometry}

\newtheorem{theorem}{Theorem}
\newtheorem{lemma}[theorem]{Lemma}
\newtheorem{definition}[theorem]{Definition}
\newtheorem{remark}[theorem]{Remark}
\newtheorem{proposition}[theorem]{Proposition}

\definecolor{proved}{RGB}{40,120,40}
\definecolor{open}{RGB}{160,50,50}
\definecolor{partial}{RGB}{160,120,0}

\title{\textbf{Exploration: \texttt{sin\_not\_in\_eml} via O-Minimal Theory}\\[0.4em]
\large What is already provable vs.\ what genuinely needs o-minimality}
\author{Arturo R.\ Almaguer}
\date{2026-04-21}

\begin{document}
\maketitle

\begin{abstract}
We map out the proof landscape for the sorry
\texttt{sin\_not\_in\_eml}: $\sin \notin \mathrm{EML}_k$ for any finite $k$.
We identify which steps can already be handled with Mathlib's current tools
(analyticity, Bolzano-Weierstrass, identity theorem) versus what genuinely
requires o-minimal structure theory (Wilkie 1996, not yet in Mathlib).
The key gap is a depth-$k$ induction for the zero-finiteness lemma
that cannot be closed without the o-minimality of $\mathbb{R}_{\exp}$.
\end{abstract}

%──────────────────────────────────────────────────────────────────────────────
\section{The Target Sorry}

\begin{verbatim}
theorem sin_not_in_eml (k : N) :
    forall t : EMLTree, t.depth <= k ->
      NOT (forall x : R, t.evalReal x = Real.sin x) := by
  sorry
\end{verbatim}

Equivalently: no finite EML tree computes $\sin(x)$ everywhere on $\mathbb{R}$.

%──────────────────────────────────────────────────────────────────────────────
\section{The Standard Proof Strategy}

The natural proof structure has four steps:

\begin{enumerate}
  \item \textbf{[PROVED]} $\sin$ has infinitely many zeros:
    $\sin(n\pi) = 0$ for all $n \in \mathbb{Z}$, and these are distinct.
    (\texttt{sin\_has\_infinitely\_many\_zeros}, 0 sorry.)

  \item \textbf{[PROVED under WFP]} Every well-formed EML tree is real-analytic on $(0,\infty)$.
    (\texttt{eml\_tree\_analytic}, 0 sorry, under \texttt{WellFormedPos}.)

  \item \textbf{[PROVED]} A non-zero analytic function on a compact interval has finitely many zeros.
    (\texttt{analytic\_finite\_zeros\_compact}, 0 sorry.)

  \item \textbf{[OPEN]} Combining: if $t.\mathtt{evalReal} = \sin$, derive a contradiction.
\end{enumerate}

Step 4 requires connecting Steps 1--3. The difficulty lies in Step 2: WFP
is a hypothesis we need to \emph{derive} from the assumption
$t.\mathtt{evalReal} = \sin$, not assume.

%──────────────────────────────────────────────────────────────────────────────
\section{The WFP Gap}

\subsection{Can WFP be derived from evalReal = sin?}

Suppose $t.\mathtt{evalReal}(x) = \sin(x)$ for all $x \in \mathbb{R}$.

\textbf{What we know}: $\sin(x)$ is a real-valued, real-analytic function.
The outer-most node of $t$ computes
\[
  t.\mathtt{evalReal}(x) = \bigl(\exp(t_1.\mathtt{eval}(\uparrow\!x)) -
  \log(t_2.\mathtt{eval}(\uparrow\!x))\bigr).\mathtt{re} = \sin(x).
\]
\textbf{What we cannot easily extract}: Whether $t_2.\mathtt{evalReal}(x) > 0$ for all $x > 0$.

The subtree $t_2$ computes some intermediate expression. From $\sin(x)$ alone, we
cannot determine the sign of $t_2$'s output without knowing the tree structure in detail.
There is no general principle saying ``if the composed result is real-valued and
equals sin, then all log arguments are positive.''

\textbf{Conclusion}: WFP cannot be derived from the equality assumption alone without
additional case analysis on the tree structure (which is exactly what o-minimality
avoids needing).

%──────────────────────────────────────────────────────────────────────────────
\section{What IS Already Provable in Lean}

\begin{proposition}[\textcolor{proved}{Provable now}]\label{prop:depth1}
$\sin \notin \mathrm{EML}_1(\mathbb{R})$: no depth-$\leq 1$ EML tree computes $\sin$.
\end{proposition}

\begin{proof}[Proof sketch]
At depth $\leq 1$, the tree $t$ has one of:
\begin{itemize}[noitemsep]
  \item \texttt{const c}: $t.\mathtt{evalReal}(x) = c.\mathtt{re}$ is constant; $\sin$ is not.
  \item \texttt{var}: $t.\mathtt{evalReal}(x) = x$; $\sin(x) \neq x$ for large $x$.
  \item $\mathtt{ceml}(\mathtt{var}, \mathtt{var})$: gives $e^x - \ln x$, not $\sin$.
  \item $\mathtt{ceml}(\mathtt{var}, \mathtt{const}\ c)$: gives $e^x - c$, not $\sin$.
  \item $\mathtt{ceml}(\mathtt{const}\ c, \mathtt{var})$: gives $c - \ln x$, not $\sin$.
  \item $\mathtt{ceml}(\mathtt{const}\ c_1, \mathtt{const}\ c_2)$: constant, not $\sin$.
\end{itemize}
In each case, the function is analytic on $(0,\infty)$ (proved by
\texttt{eml\_tree\_analytic\_depth\_le\_1}) and has finitely many zeros
(\texttt{depth1\_finite\_zeros\_real}). Since $\sin$ has infinitely many zeros,
equality is impossible. \qed
\end{proof}

\textbf{Status}: This proof is fully formalizable in Lean with existing tools.
It relies only on \texttt{eml\_tree\_analytic\_depth\_le\_1} (0 sorry) +
\texttt{analytic\_finite\_zeros\_compact} (0 sorry) + \texttt{sin\_has\_infinitely\_many\_zeros} (0 sorry).
The only missing piece is a small Lean proof connecting these three results for each depth-1 case.

\begin{proposition}[\textcolor{proved}{Provable now}]\label{prop:wfp}
If $t$ satisfies WFP on $(a,b) \subset (0,\infty)$ and $t.\mathtt{evalReal} \neq 0$
on $(a,b)$, then $t.\mathtt{evalReal}$ has at most finitely many zeros on $[a,b]$.
\end{proposition}

\begin{proof}
By \texttt{eml\_tree\_analytic} (under WFP) + \texttt{analytic\_finite\_zeros\_compact}. \qed
\end{proof}

%──────────────────────────────────────────────────────────────────────────────
\section{The O-Minimal Approach}

\subsection{O-minimality of $\mathbb{R}_{\exp}$}

\begin{theorem}[Wilkie, 1996]
The structure $(\mathbb{R},\ +,\ \cdot,\ \exp,\ <,\ 0,\ 1)$ is o-minimal.
\end{theorem}

An o-minimal structure satisfies: every definable subset of $\mathbb{R}$ is a
finite union of points and open intervals.

\subsection{Consequence for EML trees}

Every EML tree function $f = t.\mathtt{evalReal}$ is \emph{definable} in
$\mathbb{R}_{\exp}$: it is built from $\exp$ and $\ln$ (which are $\mathbb{R}_{\exp}$-definable)
by arithmetic operations. Therefore:

\begin{corollary}[from Wilkie]
For any EML tree $t$, the zero set $\{x \in \mathbb{R} \mid t.\mathtt{evalReal}(x) = 0\}$
is a finite union of points and intervals. In particular, on any bounded interval,
there are at most finitely many zeros.
\end{corollary}

This is exactly what \texttt{sin\_not\_in\_eml} needs for depth-$k$ trees.
Combined with $\sin$ having infinitely many zeros, the proof is immediate.

\subsection{What is NOT in Mathlib}

The key gap is:
\begin{itemize}[noitemsep]
  \item Wilkie's theorem is not formalized in Mathlib 4.
  \item There is no general notion of ``definable in $\mathbb{R}_{\exp}$'' in Mathlib.
  \item The o-minimal structure theory library does not exist for Lean 4.
\end{itemize}

Formalizing Wilkie's theorem is a major research project (estimated years of effort).

%──────────────────────────────────────────────────────────────────────────────
\section{Alternative Paths (not o-minimal)}

\subsection{Path A: Induction on depth using WFP as extra hypothesis}

Formulate a strengthened induction:
\begin{verbatim}
theorem sin_not_in_eml_wfp (k : N) :
    forall t : EMLTree, t.depth <= k ->
    (forall x > 0, t.WellFormedPos x) ->
    NOT (forall x : R, t.evalReal x = Real.sin x)
\end{verbatim}

Under WFP, the proof goes through: analyticity + finitely many zeros + sin's infinity.
\textbf{The remaining gap}: prove the original (without WFP hypothesis) from this.
Requires showing: if any tree computes $\sin$, it must satisfy WFP. This is unclear.

\subsection{Path B: Khovanskii's fewnomials (analytic geometry)}

Khovanskii (1991) gives explicit upper bounds on the number of zeros of \emph{Pfaffian functions}.
EML trees are Pfaffian functions (exp and log are solutions of algebraic differential equations).
The Khovanskii bound gives: an EML tree of depth $k$ has at most $B(k)$ zeros on any bounded interval,
where $B(k)$ is a computable function.

This is weaker than full o-minimality and might be more formalizable, but still requires
significant differential algebra machinery not in Mathlib.

\subsection{Path C: Direct depth-by-depth case analysis}

For each fixed $k$, prove $\sin \notin \mathrm{EML}_k$ by case analysis. Already done for $k \leq 1$
(Proposition \ref{prop:depth1}). Depth 2 is feasible but involves many cases.
This approach does not generalize to arbitrary $k$ without o-minimality.

%──────────────────────────────────────────────────────────────────────────────
\section{Summary: Provable Now vs.\ Needs O-minimal}

\begin{center}
\begin{tabular}{lll}
\toprule
Statement & Status & Tool needed \\
\midrule
$\sin$ has infinitely many zeros & \textcolor{proved}{Proved} & Lean arithmetic \\
WFP trees analytic (T\_EML\_WFP\_ANALYTIC) & \textcolor{proved}{Proved} & Mathlib analyticity \\
Analytic $\Rightarrow$ finitely many zeros & \textcolor{proved}{Proved} & Bolzano + identity thm \\
$\sin \notin \mathrm{EML}_1$ & \textcolor{proved}{Provable now} & Depth-1 case analysis \\
$\sin \notin \mathrm{EML}_k$ (fixed $k$) & \textcolor{partial}{Feasible} & Increasing case analysis \\
$\sin \notin \mathrm{EML}_k$ (all $k$) & \textcolor{open}{Open} & O-minimal theory \\
Khovanskii bound for EML trees & \textcolor{open}{Open} & Differential algebra \\
Wilkie's theorem in Lean & \textcolor{open}{Long-term} & Major project \\
\bottomrule
\end{tabular}
\end{center}

%──────────────────────────────────────────────────────────────────────────────
\section{Recommended Next Step}

The highest-leverage short-term action is \textbf{Path C for depth 2}: prove
$\sin \notin \mathrm{EML}_2$ by Lean case analysis. This:
\begin{itemize}[noitemsep]
  \item Exercises the induction infrastructure
  \item Produces a publishable partial result (together with depth 1)
  \item Does not depend on o-minimality
  \item Can be formalized in Lean 4 with current Mathlib tools
\end{itemize}

The full general theorem ($\sin \notin \mathrm{EML}_k$ for all $k$) remains genuinely
open pending Mathlib o-minimality formalization.

\end{document}
